\section{Motivation}

\subsection{The decomplected venue}

The DEX is decomplected into orthogonal homes with one concern each.
The D-Chain is the venue: it owns the per-account ledger, the order
book, the matcher, and settlement. The public-network DEX VM
(\texttt{chains/dexvm}) is a \emph{stateless atomic ZAP proxy}. The
proxy holds no order, pool, position, or balance keys; its only jobs
are (i) moving value into and out of the venue atomically across
chains via shared memory, and (ii) relaying order traffic to the
venue. The venue is the single source of truth for money and matches;
the proxy is transport. The atomic value movement reuses the cross-chain
all-or-nothing discipline of LP-211~\cite{lp211}; the relay rail rides the
ZAP universal wire of LP-201~\cite{lp201}; the proxy block is a DAG vertex in
the mempool of LP-208~\cite{lp208}.

The custody model, in one line per phase, is:

\begin{itemize}[leftmargin=2em,nosep]
\item \textbf{Deposit} $=$ \texttt{ImportTx} (an atomic debit of a
  C-chain or X-chain UTXO, consumed exactly once) \emph{then} credit
  of the D-Chain account's \texttt{available} balance.
\item \textbf{Order} $=$ move from \texttt{available} to
  \texttt{locked} (buy locks $\lceil \mathrm{size}\cdot\mathrm{price}
  \rceil$ quote units; sell locks $\mathrm{size}$ base units).
\item \textbf{Match} $=$ settle \texttt{locked} against
  \texttt{available} on \emph{both} legs, wei-exact, plus fees; every
  unit leaving one account's \texttt{locked} arrives in a
  counterparty's \texttt{available}.
\item \textbf{Withdraw} $=$ debit the D-Chain account's
  \texttt{available} \emph{then} \texttt{ExportTx} (an atomic release
  back to the originating chain).
\end{itemize}

\subsection{Why per-validator deposit/withdraw relay would fork}

The carried-fills work established the governing fact for order
submission: \textbf{a moving venue ledger makes per-validator relay
non-deterministic}. If each validator independently dialed the venue
during \texttt{Block.Verify} to relay a \texttt{clob\_submit}, each
would observe the venue ledger at a different instant and derive a
different fill set, forking the state root. The cure was to make the
\emph{proposer} relay the submit once at \texttt{BuildBlock} and carry
the resulting fills in the block and DAG-vertex bytes; every validator
then re-derives the same state transition as a pure function of
(height, carried timestamp, transaction bytes, carried fills), so all
validators reach a byte-identical state root and zero validators dial
the venue during verification.

The deposit and withdraw flows have the \emph{same} hazard for the
\emph{same} reason. Both mutate the venue ledger, and both depend on
the venue's response to a relay:

\begin{itemize}[leftmargin=2em,nosep]
\item \textbf{Deposit fork.} \texttt{executeDeposit} performs the
  atomic import (consume-once) then relays \texttt{clob\_deposit} to
  the venue, which credits \texttt{available} and refuses if the
  credited amount does not equal the imported amount. The credited
  amount, the resulting account state, and the deposit acknowledgement
  are venue responses. If each validator relayed independently during
  verification, each could observe a different ledger snapshot and
  apply a different credit (or observe a different idempotency
  outcome). The deposit acknowledgement is exactly the kind of
  ``venue said this happened'' datum that the carried-fills argument
  forbids re-deriving per validator.
\item \textbf{Withdraw fork.} \texttt{executeWithdraw} relays
  \texttt{clob\_withdraw}, which debits \texttt{available} and returns
  the \emph{realized} amount clamped to the actually-available balance,
  then performs the atomic export of exactly that realized amount. The
  realized amount is a function of the venue ledger at relay time: two
  validators relaying independently against a moving ledger could
  compute different realized amounts, hence different export
  commitments, hence different state roots --- and worse, different
  on-chain releases.
\end{itemize}

\subsection{The consequence the carried-fills work avoided, restated}

For order submission the divergence was a forked state root. For
custody the divergence is strictly worse: a withdraw realized amount
that differs across validators is a divergence in \emph{how much value
leaves the system on the originating chain}. The deposit acknowledgement
that differs across validators is a divergence in \emph{how much value
the venue believes entered}. A consensus that cannot agree on these
numbers cannot agree on the conservation invariant of
Section~\ref{sec:conservation}, and the venue is not safe to
decentralize.

\subsection{The fix, by analogy}

This LP applies the carried-fills mechanism to custody. The proposer
performs the irreversible, non-deterministic venue interaction
\emph{once} at \texttt{BuildBlock} --- the import-then-credit for
deposits, the debit-then-export-commitment for withdrawals --- and
carries the \emph{results} in the block and DAG-vertex bytes. Every
validator then applies the carried custody delta as a pure,
deterministic function of the block bytes: no validator dials the
venue, no validator re-samples a moving ledger, and the block's effect
on the state root is identical network-wide. The venue ledger remains
the source of truth; the proxy block becomes a faithful, replayable
record of the custody transitions the proposer already executed
against it. The pattern is the prepare-then-commit shape of two-phase
commit~\cite{gray1978}, specialized so the ``prepare'' is a single
proposer relay and the ``commit'' is the deterministic application of
carried results by every validator.

This is deliberately the smallest change that closes the fork: it
reuses the existing carried-payload framing, the existing
consume-once/idempotency machinery, and the reserved signature field
already present for the trustless path. It introduces no new venue
semantics --- only a deterministic way for every validator to agree on
the custody transitions the venue performed.
